\documentclass[aspectratio=169,10pt]{beamer}
\makeatletter\def\input@path{{../../estilo/}{../estilo/}{estilo/}}\makeatother
\usepackage{devmobile}
\modulo[Flutter]{12}{Firebase com FlutterFire}{Configuração, Authentication, Cloud Firestore, regras de segurança, Storage, Cloud Functions, mensagens push, Crashlytics e Emulator Suite}

\begin{document}
\capa

\begin{frame}{Objetivos do módulo}
  \begin{conceito}[Ao final deste módulo você será capaz de]
    \begin{itemize}
      \item Conectar um app Flutter ao Firebase com o FlutterFire CLI.
      \item Autenticar usuários por e-mail e senha e com Google.
      \item Modelar, ler e gravar dados no Cloud Firestore em tempo real.
      \item Proteger os dados com regras de segurança.
      \item Enviar arquivos ao Storage, chamar Cloud Functions e receber notificações push.
      \item Monitorar falhas com Crashlytics e desenvolver com o Emulator Suite.
    \end{itemize}
  \end{conceito}
\end{frame}

\sumario

% =============================================================================
%  Fundamentos compartilhados --- Módulo 12: identidade, tempo real e mensagens
%  Incluído pelas trilhas Kotlin e Flutter com % =============================================================================
%  Fundamentos compartilhados --- Módulo 12: identidade, tempo real e mensagens
%  Incluído pelas trilhas Kotlin e Flutter com % =============================================================================
%  Fundamentos compartilhados --- Módulo 12: identidade, tempo real e mensagens
%  Incluído pelas trilhas Kotlin e Flutter com \input{../fundamentos/f12-identidade-nuvem}
%  (autenticação no servidor, hash de senhas e JWT estão no Módulo 17;
%   serverless e custos, no Módulo 16; modelagem em documentos, no Módulo 09)
% =============================================================================
\section{Fundamentos: identidade, tempo real e mensagens}

\begin{frame}{Identidade federada: OAuth 2.0 e OpenID Connect}
  \begin{columns}[c]
    \begin{column}{0.5\textwidth}
      \centering
      \begin{tikzpicture}[every node/.style={font=\tiny}]
        \node[caixa,font=\tiny,minimum width=1.5cm] (app) at (0,0) {app};
        \node[caixaambar,font=\tiny,minimum width=2cm] (nav) at (3,0) {navegador do sistema};
        \node[caixaescura,font=\tiny,minimum width=2.4cm] (as) at (3,-1.4) {servidor de autorização\\ (Google, Microsoft)};
        \node[caixapetroleo,font=\tiny,minimum width=1.9cm] (api) at (0,-2.8) {API (recurso)};
        \draw[seta] (app) -- node[rotulo,above]{1. abre} (nav);
        \draw[seta] (nav) -- node[rotulo,right]{2. login e consentimento} (as);
        \draw[setaambar] (as.west) -- node[rotulo,below,sloped]{3. código $\rightarrow$ \emph{tokens}} (app.south east);
        \draw[seta] (app) -- node[rotulo,left]{4. \emph{token} de acesso} (api);
      \end{tikzpicture}
      \begin{historia}
        OAuth 2.0 virou padrão em 2012 (RFC 6749); o OpenID Connect veio em 2014. Para apps nativos, a RFC 8252 (2017) recomenda o navegador do sistema e PKCE.
      \end{historia}
    \end{column}
    \begin{column}{0.48\textwidth}
      \begin{definicao}[OAuth 2.0 e OpenID Connect]
        \textbf{OAuth 2.0} delega \textbf{autorização}: o app recebe um \emph{token} para agir em nome do usuário sem ver a senha. O \textbf{OpenID Connect} acrescenta \textbf{identidade}: um \emph{ID token} diz quem é o usuário.
      \end{definicao}
      \begin{pordentro}[PKCE]
        O app cria um segredo aleatório por login e envia só o seu \emph{hash}; na troca do código pelos \emph{tokens}, mostra o segredo. Quem interceptar o código não consegue usá-lo.
      \end{pordentro}
      \begin{atencao}
        Nunca peça a senha do Google numa tela do app nem use \emph{WebView} para o login.
      \end{atencao}
    \end{column}
  \end{columns}
\end{frame}

\begin{frame}{\emph{Tokens} no app: identidade, acesso e renovação}
  \scriptsize
  \begin{tabularx}{\textwidth}{@{}l L l L@{}}
    \toprule
    \cabecalho{\emph{Token}} & \cabecalho{Para que serve} & \cabecalho{Validade típica} & \cabecalho{Cuidados} \\
    \midrule
    ID \emph{token} & informa quem é o usuário (nome, e-mail, \cod{sub}) & 1 hora no Firebase & não serve para chamar APIs de terceiros \\
    De acesso & autoriza chamadas a uma API & minutos a 1 hora & enviar só por HTTPS, no cabeçalho \cod{Authorization} \\
    De renovação & obtém novos \emph{tokens} sem novo login & semanas a meses & guardar em armazenamento seguro; revogar no \emph{logout} \\
    \bottomrule
  \end{tabularx}
  \normalsize
  \vspace{4pt}
  \begin{columns}[T]
    \begin{column}{0.49\textwidth}
      \begin{conceito}[Vida curta, renovação automática]
        \emph{Tokens} de acesso curtos limitam o estrago se vazarem. Os SDKs do Firebase renovam o ID \emph{token} sozinhos antes de expirar.
      \end{conceito}
    \end{column}
    \begin{column}{0.49\textwidth}
      \begin{pordentro}[Papéis dentro do \emph{token}]
        O servidor pode gravar papéis no \emph{token} (\emph{custom claims}, como \cod{professor: true}). As regras de segurança e as funções leem esses papéis sem consultar o banco.
      \end{pordentro}
    \end{column}
  \end{columns}
\end{frame}

\begin{frame}{Senhas, múltiplos fatores e \emph{passkeys}}
  \begin{columns}[T]
    \begin{column}{0.5\textwidth}
      \begin{definicao}[Fatores de autenticação]
        Algo que você \textbf{sabe} (senha), algo que você \textbf{tem} (celular, chave física) e algo que você \textbf{é} (digital, rosto). Combinar dois fatores diferentes é a autenticação multifator (MFA).
      \end{definicao}
      \begin{atencao}[O ponto fraco da senha]
        Senhas são reutilizadas, vazam e podem ser digitadas num site falso (\emph{phishing}). Códigos por SMS ajudam, mas também podem ser interceptados.
      \end{atencao}
    \end{column}
    \begin{column}{0.48\textwidth}
      \begin{definicao}[\emph{Passkey}]
        Credencial baseada em par de chaves, criada para um único site ou app. A chave privada fica no aparelho ou no gerenciador de senhas; o servidor guarda só a pública.
      \end{definicao}
      \begin{pordentro}[Por que resiste a \emph{phishing}]
        O login é a assinatura de um desafio do servidor, liberada pela biometria ou pelo PIN do aparelho. A chave só funciona no domínio para o qual foi criada, e não há segredo para vazar do servidor.
      \end{pordentro}
    \end{column}
  \end{columns}
\end{frame}

\begin{frame}{Bancos em tempo real: escuta, cache local e compensação de latência}
  \begin{columns}[c]
    \begin{column}{0.48\textwidth}
      \centering
      \begin{tikzpicture}[every node/.style={font=\tiny}]
        \node[caixa,font=\tiny,minimum width=1.9cm] (ui) at (0,0) {tela};
        \node[caixaambar,font=\tiny,minimum width=1.9cm] (cache) at (0,-1.3) {cache local do SDK};
        \node[caixaescura,font=\tiny,minimum width=1.9cm] (srv) at (3.6,-1.3) {banco na nuvem};
        \node[caixa,font=\tiny,minimum width=1.9cm] (outro) at (3.6,0) {outro aparelho};
        \draw[setaambar] ([xshift=-4mm]ui.south) -- node[rotulo,left]{1. grava} ([xshift=-4mm]cache.north);
        \draw[seta] ([xshift=4mm]cache.north) -- node[rotulo,right,text width=1.3cm]{2. avisa na hora (pendente)} ([xshift=4mm]ui.south);
        \draw[seta] ([yshift=1mm]cache.east) -- node[rotulo,above]{3. envia} ([yshift=1mm]srv.west);
        \draw[seta] ([yshift=-1mm]srv.west) -- node[rotulo,below]{4. confirma} ([yshift=-1mm]cache.east);
        \draw[seta] (srv) -- node[rotulo,right]{5. notifica} (outro);
      \end{tikzpicture}
    \end{column}
    \begin{column}{0.5\textwidth}
      \begin{definicao}[Escuta (\emph{listener})]
        Em vez de perguntar repetidas vezes, o app assina uma consulta e recebe o resultado inicial e depois cada mudança, por uma conexão mantida aberta.
      \end{definicao}
      \begin{definicao}[Compensação de latência]
        A gravação aparece na tela antes de o servidor confirmar, marcada como pendente. Se o servidor recusar, o SDK desfaz a mudança.
      \end{definicao}
      \begin{atencao}[Escutas custam]
        Cada documento recebido conta como leitura. Limite as consultas (filtros, \cod{limit}) e encerre as escutas quando a tela sair.
      \end{atencao}
    \end{column}
  \end{columns}
\end{frame}

\begin{frame}[fragile]{Regras de segurança: a autorização mora no servidor}
  \begin{columns}[T]
    \begin{column}{0.52\textwidth}
      \begin{conceito}[O app fala direto com o banco]
        Em um BaaS não há servidor seu no meio: qualquer pessoa pode usar as chaves públicas do app em outro programa. Verificações feitas no app são só experiência de uso; a proteção real são as \textbf{regras}, avaliadas no servidor a cada requisição.
      \end{conceito}
\begin{codigo}[title=firestore.rules (trecho)]{JSON}
match /reservas/{id} {
  allow read: if request.auth.uid == resource.data.uid;
  allow create: if request.auth.uid
                   == request.resource.data.uid;
}
\end{codigo}
    \end{column}
    \begin{column}{0.46\textwidth}
      \begin{definicao}[Menor privilégio]
        Comece negando tudo e libere só o necessário, para quem precisa. Valide também o formato: campos obrigatórios, tipos e tamanhos.
      \end{definicao}
      \begin{atencao}[Regras não são filtros]
        A regra não esconde documentos de uma consulta: se a consulta \textbf{puder} trazer algo proibido, ela inteira falha. A consulta precisa ter os mesmos filtros da regra.
      \end{atencao}
      \begin{dica}
        Teste as regras no emulador, com casos permitidos e negados, como qualquer outro código.
      \end{dica}
    \end{column}
  \end{columns}
\end{frame}

\begin{frame}{Notificações \emph{push} por dentro}
  \centering
  \begin{tikzpicture}[every node/.style={font=\tiny}]
    \node[caixa,font=\tiny,minimum width=1.8cm,minimum height=9mm] (app) at (0.9,0) {app no aparelho};
    \node[caixaescura,font=\tiny,minimum width=2cm,minimum height=9mm] (srv) at (4.1,0) {seu servidor\\ ou função};
    \node[caixaambarforte,font=\tiny\bfseries,minimum width=1.4cm,minimum height=9mm] (fcm) at (7.1,0) {FCM};
    \node[caixapetroleo,font=\tiny,minimum width=1.6cm,minimum height=9mm] (apns) at (9.9,0) {APNs\\ (iPhone)};
    \node[caixa,font=\tiny,minimum width=2.2cm,minimum height=9mm] (dev) at (13.1,0) {aparelho: conexão\\ única do sistema};
    \draw[seta] (app) -- node[rotulo,above]{1. \emph{token}} (srv);
    \draw[seta] (srv) -- node[rotulo,above]{2. HTTP v1} (fcm);
    \draw[seta] (fcm) -- node[rotulo,above]{3. iOS} (apns);
    \draw[seta] (apns) -- node[rotulo,above]{4. entrega} (dev);
    \draw[setaambar,rounded corners] (fcm.south) -- ++(0,-4mm) -| node[rotulo,pos=0.25,below]{3. Android: direto, pelo Google Play services} (dev.south);
\end{tikzpicture}
  \raggedright
  \begin{columns}[T]
    \begin{column}{0.32\textwidth}
      \begin{definicao}[\emph{Token} de registro]
        Endereço do app naquele aparelho. Muda em reinstalações e após limpar dados: mantenha-o atualizado no servidor.
      \end{definicao}
    \end{column}
    \begin{column}{0.32\textwidth}
      \begin{conceito}[Dois tipos de mensagem]
        \textbf{Notificação}: o sistema exibe sozinho. \textbf{Dados}: o código do app decide o que fazer.
      \end{conceito}
    \end{column}
    \begin{column}{0.32\textwidth}
      \begin{atencao}[Entrega não garantida]
        O usuário pode negar a permissão (Android 13+ e iOS), e o sistema adia mensagens para poupar bateria.
      \end{atencao}
    \end{column}
  \end{columns}
\end{frame}

\begin{frame}{Funções acionadas por eventos}
  \begin{columns}[c]
    \begin{column}{0.46\textwidth}
      \centering
      \begin{tikzpicture}[node distance=4mm,every node/.style={font=\tiny}]
        \node[caixa,font=\tiny,minimum width=3.6cm] (e) {evento: nova leitura gravada};
        \node[caixaescura,font=\tiny,minimum width=3.6cm,below=of e] (f) {função: valor acima do limite?};
        \node[caixaambar,font=\tiny,minimum width=1.7cm,below left=4mm and -17mm of f] (p) {envia \emph{push}};
        \node[caixaambar,font=\tiny,minimum width=1.7cm,below right=4mm and -17mm of f] (a) {atualiza resumo};
        \draw[seta] (e) -- (f); \draw[seta] (f) -- (p); \draw[seta] (f) -- (a);
      \end{tikzpicture}
      \begin{dica}
        Outros gatilhos: requisição HTTP, conta criada, arquivo enviado e horário agendado.
      \end{dica}
    \end{column}
    \begin{column}{0.52\textwidth}
      \begin{definicao}[Arquitetura orientada a eventos]
        Componentes reagem a fatos (``uma leitura foi gravada'') em vez de chamarem uns aos outros. Quem grava não precisa saber quem reage.
      \end{definicao}
      \begin{atencao}[Pelo menos uma vez]
        O provedor garante executar a função ao menos uma vez por evento, e às vezes executa duas. Escreva funções \textbf{idempotentes}: processar o mesmo evento de novo não pode duplicar o \emph{push} nem o resumo.
      \end{atencao}
      \begin{pordentro}[Onde guardar segredos]
        Chaves de serviços pagos ficam no servidor e nunca no app. O App Check (Play Integrity, App Attest) ajuda a recusar chamadas que não vêm do app legítimo.
      \end{pordentro}
    \end{column}
  \end{columns}
\end{frame}

%  (autenticação no servidor, hash de senhas e JWT estão no Módulo 17;
%   serverless e custos, no Módulo 16; modelagem em documentos, no Módulo 09)
% =============================================================================
\section{Fundamentos: identidade, tempo real e mensagens}

\begin{frame}{Identidade federada: OAuth 2.0 e OpenID Connect}
  \begin{columns}[c]
    \begin{column}{0.5\textwidth}
      \centering
      \begin{tikzpicture}[every node/.style={font=\tiny}]
        \node[caixa,font=\tiny,minimum width=1.5cm] (app) at (0,0) {app};
        \node[caixaambar,font=\tiny,minimum width=2cm] (nav) at (3,0) {navegador do sistema};
        \node[caixaescura,font=\tiny,minimum width=2.4cm] (as) at (3,-1.4) {servidor de autorização\\ (Google, Microsoft)};
        \node[caixapetroleo,font=\tiny,minimum width=1.9cm] (api) at (0,-2.8) {API (recurso)};
        \draw[seta] (app) -- node[rotulo,above]{1. abre} (nav);
        \draw[seta] (nav) -- node[rotulo,right]{2. login e consentimento} (as);
        \draw[setaambar] (as.west) -- node[rotulo,below,sloped]{3. código $\rightarrow$ \emph{tokens}} (app.south east);
        \draw[seta] (app) -- node[rotulo,left]{4. \emph{token} de acesso} (api);
      \end{tikzpicture}
      \begin{historia}
        OAuth 2.0 virou padrão em 2012 (RFC 6749); o OpenID Connect veio em 2014. Para apps nativos, a RFC 8252 (2017) recomenda o navegador do sistema e PKCE.
      \end{historia}
    \end{column}
    \begin{column}{0.48\textwidth}
      \begin{definicao}[OAuth 2.0 e OpenID Connect]
        \textbf{OAuth 2.0} delega \textbf{autorização}: o app recebe um \emph{token} para agir em nome do usuário sem ver a senha. O \textbf{OpenID Connect} acrescenta \textbf{identidade}: um \emph{ID token} diz quem é o usuário.
      \end{definicao}
      \begin{pordentro}[PKCE]
        O app cria um segredo aleatório por login e envia só o seu \emph{hash}; na troca do código pelos \emph{tokens}, mostra o segredo. Quem interceptar o código não consegue usá-lo.
      \end{pordentro}
      \begin{atencao}
        Nunca peça a senha do Google numa tela do app nem use \emph{WebView} para o login.
      \end{atencao}
    \end{column}
  \end{columns}
\end{frame}

\begin{frame}{\emph{Tokens} no app: identidade, acesso e renovação}
  \scriptsize
  \begin{tabularx}{\textwidth}{@{}l L l L@{}}
    \toprule
    \cabecalho{\emph{Token}} & \cabecalho{Para que serve} & \cabecalho{Validade típica} & \cabecalho{Cuidados} \\
    \midrule
    ID \emph{token} & informa quem é o usuário (nome, e-mail, \cod{sub}) & 1 hora no Firebase & não serve para chamar APIs de terceiros \\
    De acesso & autoriza chamadas a uma API & minutos a 1 hora & enviar só por HTTPS, no cabeçalho \cod{Authorization} \\
    De renovação & obtém novos \emph{tokens} sem novo login & semanas a meses & guardar em armazenamento seguro; revogar no \emph{logout} \\
    \bottomrule
  \end{tabularx}
  \normalsize
  \vspace{4pt}
  \begin{columns}[T]
    \begin{column}{0.49\textwidth}
      \begin{conceito}[Vida curta, renovação automática]
        \emph{Tokens} de acesso curtos limitam o estrago se vazarem. Os SDKs do Firebase renovam o ID \emph{token} sozinhos antes de expirar.
      \end{conceito}
    \end{column}
    \begin{column}{0.49\textwidth}
      \begin{pordentro}[Papéis dentro do \emph{token}]
        O servidor pode gravar papéis no \emph{token} (\emph{custom claims}, como \cod{professor: true}). As regras de segurança e as funções leem esses papéis sem consultar o banco.
      \end{pordentro}
    \end{column}
  \end{columns}
\end{frame}

\begin{frame}{Senhas, múltiplos fatores e \emph{passkeys}}
  \begin{columns}[T]
    \begin{column}{0.5\textwidth}
      \begin{definicao}[Fatores de autenticação]
        Algo que você \textbf{sabe} (senha), algo que você \textbf{tem} (celular, chave física) e algo que você \textbf{é} (digital, rosto). Combinar dois fatores diferentes é a autenticação multifator (MFA).
      \end{definicao}
      \begin{atencao}[O ponto fraco da senha]
        Senhas são reutilizadas, vazam e podem ser digitadas num site falso (\emph{phishing}). Códigos por SMS ajudam, mas também podem ser interceptados.
      \end{atencao}
    \end{column}
    \begin{column}{0.48\textwidth}
      \begin{definicao}[\emph{Passkey}]
        Credencial baseada em par de chaves, criada para um único site ou app. A chave privada fica no aparelho ou no gerenciador de senhas; o servidor guarda só a pública.
      \end{definicao}
      \begin{pordentro}[Por que resiste a \emph{phishing}]
        O login é a assinatura de um desafio do servidor, liberada pela biometria ou pelo PIN do aparelho. A chave só funciona no domínio para o qual foi criada, e não há segredo para vazar do servidor.
      \end{pordentro}
    \end{column}
  \end{columns}
\end{frame}

\begin{frame}{Bancos em tempo real: escuta, cache local e compensação de latência}
  \begin{columns}[c]
    \begin{column}{0.48\textwidth}
      \centering
      \begin{tikzpicture}[every node/.style={font=\tiny}]
        \node[caixa,font=\tiny,minimum width=1.9cm] (ui) at (0,0) {tela};
        \node[caixaambar,font=\tiny,minimum width=1.9cm] (cache) at (0,-1.3) {cache local do SDK};
        \node[caixaescura,font=\tiny,minimum width=1.9cm] (srv) at (3.6,-1.3) {banco na nuvem};
        \node[caixa,font=\tiny,minimum width=1.9cm] (outro) at (3.6,0) {outro aparelho};
        \draw[setaambar] ([xshift=-4mm]ui.south) -- node[rotulo,left]{1. grava} ([xshift=-4mm]cache.north);
        \draw[seta] ([xshift=4mm]cache.north) -- node[rotulo,right,text width=1.3cm]{2. avisa na hora (pendente)} ([xshift=4mm]ui.south);
        \draw[seta] ([yshift=1mm]cache.east) -- node[rotulo,above]{3. envia} ([yshift=1mm]srv.west);
        \draw[seta] ([yshift=-1mm]srv.west) -- node[rotulo,below]{4. confirma} ([yshift=-1mm]cache.east);
        \draw[seta] (srv) -- node[rotulo,right]{5. notifica} (outro);
      \end{tikzpicture}
    \end{column}
    \begin{column}{0.5\textwidth}
      \begin{definicao}[Escuta (\emph{listener})]
        Em vez de perguntar repetidas vezes, o app assina uma consulta e recebe o resultado inicial e depois cada mudança, por uma conexão mantida aberta.
      \end{definicao}
      \begin{definicao}[Compensação de latência]
        A gravação aparece na tela antes de o servidor confirmar, marcada como pendente. Se o servidor recusar, o SDK desfaz a mudança.
      \end{definicao}
      \begin{atencao}[Escutas custam]
        Cada documento recebido conta como leitura. Limite as consultas (filtros, \cod{limit}) e encerre as escutas quando a tela sair.
      \end{atencao}
    \end{column}
  \end{columns}
\end{frame}

\begin{frame}[fragile]{Regras de segurança: a autorização mora no servidor}
  \begin{columns}[T]
    \begin{column}{0.52\textwidth}
      \begin{conceito}[O app fala direto com o banco]
        Em um BaaS não há servidor seu no meio: qualquer pessoa pode usar as chaves públicas do app em outro programa. Verificações feitas no app são só experiência de uso; a proteção real são as \textbf{regras}, avaliadas no servidor a cada requisição.
      \end{conceito}
\begin{codigo}[title=firestore.rules (trecho)]{JSON}
match /reservas/{id} {
  allow read: if request.auth.uid == resource.data.uid;
  allow create: if request.auth.uid
                   == request.resource.data.uid;
}
\end{codigo}
    \end{column}
    \begin{column}{0.46\textwidth}
      \begin{definicao}[Menor privilégio]
        Comece negando tudo e libere só o necessário, para quem precisa. Valide também o formato: campos obrigatórios, tipos e tamanhos.
      \end{definicao}
      \begin{atencao}[Regras não são filtros]
        A regra não esconde documentos de uma consulta: se a consulta \textbf{puder} trazer algo proibido, ela inteira falha. A consulta precisa ter os mesmos filtros da regra.
      \end{atencao}
      \begin{dica}
        Teste as regras no emulador, com casos permitidos e negados, como qualquer outro código.
      \end{dica}
    \end{column}
  \end{columns}
\end{frame}

\begin{frame}{Notificações \emph{push} por dentro}
  \centering
  \begin{tikzpicture}[every node/.style={font=\tiny}]
    \node[caixa,font=\tiny,minimum width=1.8cm,minimum height=9mm] (app) at (0.9,0) {app no aparelho};
    \node[caixaescura,font=\tiny,minimum width=2cm,minimum height=9mm] (srv) at (4.1,0) {seu servidor\\ ou função};
    \node[caixaambarforte,font=\tiny\bfseries,minimum width=1.4cm,minimum height=9mm] (fcm) at (7.1,0) {FCM};
    \node[caixapetroleo,font=\tiny,minimum width=1.6cm,minimum height=9mm] (apns) at (9.9,0) {APNs\\ (iPhone)};
    \node[caixa,font=\tiny,minimum width=2.2cm,minimum height=9mm] (dev) at (13.1,0) {aparelho: conexão\\ única do sistema};
    \draw[seta] (app) -- node[rotulo,above]{1. \emph{token}} (srv);
    \draw[seta] (srv) -- node[rotulo,above]{2. HTTP v1} (fcm);
    \draw[seta] (fcm) -- node[rotulo,above]{3. iOS} (apns);
    \draw[seta] (apns) -- node[rotulo,above]{4. entrega} (dev);
    \draw[setaambar,rounded corners] (fcm.south) -- ++(0,-4mm) -| node[rotulo,pos=0.25,below]{3. Android: direto, pelo Google Play services} (dev.south);
\end{tikzpicture}
  \raggedright
  \begin{columns}[T]
    \begin{column}{0.32\textwidth}
      \begin{definicao}[\emph{Token} de registro]
        Endereço do app naquele aparelho. Muda em reinstalações e após limpar dados: mantenha-o atualizado no servidor.
      \end{definicao}
    \end{column}
    \begin{column}{0.32\textwidth}
      \begin{conceito}[Dois tipos de mensagem]
        \textbf{Notificação}: o sistema exibe sozinho. \textbf{Dados}: o código do app decide o que fazer.
      \end{conceito}
    \end{column}
    \begin{column}{0.32\textwidth}
      \begin{atencao}[Entrega não garantida]
        O usuário pode negar a permissão (Android 13+ e iOS), e o sistema adia mensagens para poupar bateria.
      \end{atencao}
    \end{column}
  \end{columns}
\end{frame}

\begin{frame}{Funções acionadas por eventos}
  \begin{columns}[c]
    \begin{column}{0.46\textwidth}
      \centering
      \begin{tikzpicture}[node distance=4mm,every node/.style={font=\tiny}]
        \node[caixa,font=\tiny,minimum width=3.6cm] (e) {evento: nova leitura gravada};
        \node[caixaescura,font=\tiny,minimum width=3.6cm,below=of e] (f) {função: valor acima do limite?};
        \node[caixaambar,font=\tiny,minimum width=1.7cm,below left=4mm and -17mm of f] (p) {envia \emph{push}};
        \node[caixaambar,font=\tiny,minimum width=1.7cm,below right=4mm and -17mm of f] (a) {atualiza resumo};
        \draw[seta] (e) -- (f); \draw[seta] (f) -- (p); \draw[seta] (f) -- (a);
      \end{tikzpicture}
      \begin{dica}
        Outros gatilhos: requisição HTTP, conta criada, arquivo enviado e horário agendado.
      \end{dica}
    \end{column}
    \begin{column}{0.52\textwidth}
      \begin{definicao}[Arquitetura orientada a eventos]
        Componentes reagem a fatos (``uma leitura foi gravada'') em vez de chamarem uns aos outros. Quem grava não precisa saber quem reage.
      \end{definicao}
      \begin{atencao}[Pelo menos uma vez]
        O provedor garante executar a função ao menos uma vez por evento, e às vezes executa duas. Escreva funções \textbf{idempotentes}: processar o mesmo evento de novo não pode duplicar o \emph{push} nem o resumo.
      \end{atencao}
      \begin{pordentro}[Onde guardar segredos]
        Chaves de serviços pagos ficam no servidor e nunca no app. O App Check (Play Integrity, App Attest) ajuda a recusar chamadas que não vêm do app legítimo.
      \end{pordentro}
    \end{column}
  \end{columns}
\end{frame}

%  (autenticação no servidor, hash de senhas e JWT estão no Módulo 17;
%   serverless e custos, no Módulo 16; modelagem em documentos, no Módulo 09)
% =============================================================================
\section{Fundamentos: identidade, tempo real e mensagens}

\begin{frame}{Identidade federada: OAuth 2.0 e OpenID Connect}
  \begin{columns}[c]
    \begin{column}{0.5\textwidth}
      \centering
      \begin{tikzpicture}[every node/.style={font=\tiny}]
        \node[caixa,font=\tiny,minimum width=1.5cm] (app) at (0,0) {app};
        \node[caixaambar,font=\tiny,minimum width=2cm] (nav) at (3,0) {navegador do sistema};
        \node[caixaescura,font=\tiny,minimum width=2.4cm] (as) at (3,-1.4) {servidor de autorização\\ (Google, Microsoft)};
        \node[caixapetroleo,font=\tiny,minimum width=1.9cm] (api) at (0,-2.8) {API (recurso)};
        \draw[seta] (app) -- node[rotulo,above]{1. abre} (nav);
        \draw[seta] (nav) -- node[rotulo,right]{2. login e consentimento} (as);
        \draw[setaambar] (as.west) -- node[rotulo,below,sloped]{3. código $\rightarrow$ \emph{tokens}} (app.south east);
        \draw[seta] (app) -- node[rotulo,left]{4. \emph{token} de acesso} (api);
      \end{tikzpicture}
      \begin{historia}
        OAuth 2.0 virou padrão em 2012 (RFC 6749); o OpenID Connect veio em 2014. Para apps nativos, a RFC 8252 (2017) recomenda o navegador do sistema e PKCE.
      \end{historia}
    \end{column}
    \begin{column}{0.48\textwidth}
      \begin{definicao}[OAuth 2.0 e OpenID Connect]
        \textbf{OAuth 2.0} delega \textbf{autorização}: o app recebe um \emph{token} para agir em nome do usuário sem ver a senha. O \textbf{OpenID Connect} acrescenta \textbf{identidade}: um \emph{ID token} diz quem é o usuário.
      \end{definicao}
      \begin{pordentro}[PKCE]
        O app cria um segredo aleatório por login e envia só o seu \emph{hash}; na troca do código pelos \emph{tokens}, mostra o segredo. Quem interceptar o código não consegue usá-lo.
      \end{pordentro}
      \begin{atencao}
        Nunca peça a senha do Google numa tela do app nem use \emph{WebView} para o login.
      \end{atencao}
    \end{column}
  \end{columns}
\end{frame}

\begin{frame}{\emph{Tokens} no app: identidade, acesso e renovação}
  \scriptsize
  \begin{tabularx}{\textwidth}{@{}l L l L@{}}
    \toprule
    \cabecalho{\emph{Token}} & \cabecalho{Para que serve} & \cabecalho{Validade típica} & \cabecalho{Cuidados} \\
    \midrule
    ID \emph{token} & informa quem é o usuário (nome, e-mail, \cod{sub}) & 1 hora no Firebase & não serve para chamar APIs de terceiros \\
    De acesso & autoriza chamadas a uma API & minutos a 1 hora & enviar só por HTTPS, no cabeçalho \cod{Authorization} \\
    De renovação & obtém novos \emph{tokens} sem novo login & semanas a meses & guardar em armazenamento seguro; revogar no \emph{logout} \\
    \bottomrule
  \end{tabularx}
  \normalsize
  \vspace{4pt}
  \begin{columns}[T]
    \begin{column}{0.49\textwidth}
      \begin{conceito}[Vida curta, renovação automática]
        \emph{Tokens} de acesso curtos limitam o estrago se vazarem. Os SDKs do Firebase renovam o ID \emph{token} sozinhos antes de expirar.
      \end{conceito}
    \end{column}
    \begin{column}{0.49\textwidth}
      \begin{pordentro}[Papéis dentro do \emph{token}]
        O servidor pode gravar papéis no \emph{token} (\emph{custom claims}, como \cod{professor: true}). As regras de segurança e as funções leem esses papéis sem consultar o banco.
      \end{pordentro}
    \end{column}
  \end{columns}
\end{frame}

\begin{frame}{Senhas, múltiplos fatores e \emph{passkeys}}
  \begin{columns}[T]
    \begin{column}{0.5\textwidth}
      \begin{definicao}[Fatores de autenticação]
        Algo que você \textbf{sabe} (senha), algo que você \textbf{tem} (celular, chave física) e algo que você \textbf{é} (digital, rosto). Combinar dois fatores diferentes é a autenticação multifator (MFA).
      \end{definicao}
      \begin{atencao}[O ponto fraco da senha]
        Senhas são reutilizadas, vazam e podem ser digitadas num site falso (\emph{phishing}). Códigos por SMS ajudam, mas também podem ser interceptados.
      \end{atencao}
    \end{column}
    \begin{column}{0.48\textwidth}
      \begin{definicao}[\emph{Passkey}]
        Credencial baseada em par de chaves, criada para um único site ou app. A chave privada fica no aparelho ou no gerenciador de senhas; o servidor guarda só a pública.
      \end{definicao}
      \begin{pordentro}[Por que resiste a \emph{phishing}]
        O login é a assinatura de um desafio do servidor, liberada pela biometria ou pelo PIN do aparelho. A chave só funciona no domínio para o qual foi criada, e não há segredo para vazar do servidor.
      \end{pordentro}
    \end{column}
  \end{columns}
\end{frame}

\begin{frame}{Bancos em tempo real: escuta, cache local e compensação de latência}
  \begin{columns}[c]
    \begin{column}{0.48\textwidth}
      \centering
      \begin{tikzpicture}[every node/.style={font=\tiny}]
        \node[caixa,font=\tiny,minimum width=1.9cm] (ui) at (0,0) {tela};
        \node[caixaambar,font=\tiny,minimum width=1.9cm] (cache) at (0,-1.3) {cache local do SDK};
        \node[caixaescura,font=\tiny,minimum width=1.9cm] (srv) at (3.6,-1.3) {banco na nuvem};
        \node[caixa,font=\tiny,minimum width=1.9cm] (outro) at (3.6,0) {outro aparelho};
        \draw[setaambar] ([xshift=-4mm]ui.south) -- node[rotulo,left]{1. grava} ([xshift=-4mm]cache.north);
        \draw[seta] ([xshift=4mm]cache.north) -- node[rotulo,right,text width=1.3cm]{2. avisa na hora (pendente)} ([xshift=4mm]ui.south);
        \draw[seta] ([yshift=1mm]cache.east) -- node[rotulo,above]{3. envia} ([yshift=1mm]srv.west);
        \draw[seta] ([yshift=-1mm]srv.west) -- node[rotulo,below]{4. confirma} ([yshift=-1mm]cache.east);
        \draw[seta] (srv) -- node[rotulo,right]{5. notifica} (outro);
      \end{tikzpicture}
    \end{column}
    \begin{column}{0.5\textwidth}
      \begin{definicao}[Escuta (\emph{listener})]
        Em vez de perguntar repetidas vezes, o app assina uma consulta e recebe o resultado inicial e depois cada mudança, por uma conexão mantida aberta.
      \end{definicao}
      \begin{definicao}[Compensação de latência]
        A gravação aparece na tela antes de o servidor confirmar, marcada como pendente. Se o servidor recusar, o SDK desfaz a mudança.
      \end{definicao}
      \begin{atencao}[Escutas custam]
        Cada documento recebido conta como leitura. Limite as consultas (filtros, \cod{limit}) e encerre as escutas quando a tela sair.
      \end{atencao}
    \end{column}
  \end{columns}
\end{frame}

\begin{frame}[fragile]{Regras de segurança: a autorização mora no servidor}
  \begin{columns}[T]
    \begin{column}{0.52\textwidth}
      \begin{conceito}[O app fala direto com o banco]
        Em um BaaS não há servidor seu no meio: qualquer pessoa pode usar as chaves públicas do app em outro programa. Verificações feitas no app são só experiência de uso; a proteção real são as \textbf{regras}, avaliadas no servidor a cada requisição.
      \end{conceito}
\begin{codigo}[title=firestore.rules (trecho)]{JSON}
match /reservas/{id} {
  allow read: if request.auth.uid == resource.data.uid;
  allow create: if request.auth.uid
                   == request.resource.data.uid;
}
\end{codigo}
    \end{column}
    \begin{column}{0.46\textwidth}
      \begin{definicao}[Menor privilégio]
        Comece negando tudo e libere só o necessário, para quem precisa. Valide também o formato: campos obrigatórios, tipos e tamanhos.
      \end{definicao}
      \begin{atencao}[Regras não são filtros]
        A regra não esconde documentos de uma consulta: se a consulta \textbf{puder} trazer algo proibido, ela inteira falha. A consulta precisa ter os mesmos filtros da regra.
      \end{atencao}
      \begin{dica}
        Teste as regras no emulador, com casos permitidos e negados, como qualquer outro código.
      \end{dica}
    \end{column}
  \end{columns}
\end{frame}

\begin{frame}{Notificações \emph{push} por dentro}
  \centering
  \begin{tikzpicture}[every node/.style={font=\tiny}]
    \node[caixa,font=\tiny,minimum width=1.8cm,minimum height=9mm] (app) at (0.9,0) {app no aparelho};
    \node[caixaescura,font=\tiny,minimum width=2cm,minimum height=9mm] (srv) at (4.1,0) {seu servidor\\ ou função};
    \node[caixaambarforte,font=\tiny\bfseries,minimum width=1.4cm,minimum height=9mm] (fcm) at (7.1,0) {FCM};
    \node[caixapetroleo,font=\tiny,minimum width=1.6cm,minimum height=9mm] (apns) at (9.9,0) {APNs\\ (iPhone)};
    \node[caixa,font=\tiny,minimum width=2.2cm,minimum height=9mm] (dev) at (13.1,0) {aparelho: conexão\\ única do sistema};
    \draw[seta] (app) -- node[rotulo,above]{1. \emph{token}} (srv);
    \draw[seta] (srv) -- node[rotulo,above]{2. HTTP v1} (fcm);
    \draw[seta] (fcm) -- node[rotulo,above]{3. iOS} (apns);
    \draw[seta] (apns) -- node[rotulo,above]{4. entrega} (dev);
    \draw[setaambar,rounded corners] (fcm.south) -- ++(0,-4mm) -| node[rotulo,pos=0.25,below]{3. Android: direto, pelo Google Play services} (dev.south);
\end{tikzpicture}
  \raggedright
  \begin{columns}[T]
    \begin{column}{0.32\textwidth}
      \begin{definicao}[\emph{Token} de registro]
        Endereço do app naquele aparelho. Muda em reinstalações e após limpar dados: mantenha-o atualizado no servidor.
      \end{definicao}
    \end{column}
    \begin{column}{0.32\textwidth}
      \begin{conceito}[Dois tipos de mensagem]
        \textbf{Notificação}: o sistema exibe sozinho. \textbf{Dados}: o código do app decide o que fazer.
      \end{conceito}
    \end{column}
    \begin{column}{0.32\textwidth}
      \begin{atencao}[Entrega não garantida]
        O usuário pode negar a permissão (Android 13+ e iOS), e o sistema adia mensagens para poupar bateria.
      \end{atencao}
    \end{column}
  \end{columns}
\end{frame}

\begin{frame}{Funções acionadas por eventos}
  \begin{columns}[c]
    \begin{column}{0.46\textwidth}
      \centering
      \begin{tikzpicture}[node distance=4mm,every node/.style={font=\tiny}]
        \node[caixa,font=\tiny,minimum width=3.6cm] (e) {evento: nova leitura gravada};
        \node[caixaescura,font=\tiny,minimum width=3.6cm,below=of e] (f) {função: valor acima do limite?};
        \node[caixaambar,font=\tiny,minimum width=1.7cm,below left=4mm and -17mm of f] (p) {envia \emph{push}};
        \node[caixaambar,font=\tiny,minimum width=1.7cm,below right=4mm and -17mm of f] (a) {atualiza resumo};
        \draw[seta] (e) -- (f); \draw[seta] (f) -- (p); \draw[seta] (f) -- (a);
      \end{tikzpicture}
      \begin{dica}
        Outros gatilhos: requisição HTTP, conta criada, arquivo enviado e horário agendado.
      \end{dica}
    \end{column}
    \begin{column}{0.52\textwidth}
      \begin{definicao}[Arquitetura orientada a eventos]
        Componentes reagem a fatos (``uma leitura foi gravada'') em vez de chamarem uns aos outros. Quem grava não precisa saber quem reage.
      \end{definicao}
      \begin{atencao}[Pelo menos uma vez]
        O provedor garante executar a função ao menos uma vez por evento, e às vezes executa duas. Escreva funções \textbf{idempotentes}: processar o mesmo evento de novo não pode duplicar o \emph{push} nem o resumo.
      \end{atencao}
      \begin{pordentro}[Onde guardar segredos]
        Chaves de serviços pagos ficam no servidor e nunca no app. O App Check (Play Integrity, App Attest) ajuda a recusar chamadas que não vêm do app legítimo.
      \end{pordentro}
    \end{column}
  \end{columns}
\end{frame}


% =============================================================================
\section{Visão geral}

\begin{frame}{FlutterFire: os plugins do Firebase}
  \begin{columns}[T]
    \begin{column}{0.56\textwidth}
      \scriptsize
      \begin{tabularx}{\linewidth}{@{}p{3.1cm} L@{}}
        \toprule
        \cabecalho{Pacote} & \cabecalho{Serviço} \\
        \midrule
        \cod{firebase\_core} & inicialização (obrigatório) \\
        \cod{firebase\_auth} & autenticação \\
        \cod{cloud\_firestore} & banco NoSQL em tempo real \\
        \cod{firebase\_storage} & arquivos (fotos, PDFs) \\
        \cod{cloud\_functions} & código no servidor \\
        \cod{firebase\_messaging} & notificações push (FCM) \\
        \cod{firebase\_crashlytics} & relatórios de falhas \\
        \cod{firebase\_analytics} & eventos de uso \\
        \cod{firebase\_remote\_config} & configuração remota \\
        \cod{firebase\_app\_check} & só o seu app acessa o backend \\
        \cod{firebase\_ai} & modelos Gemini (Módulo 20) \\
        \bottomrule
      \end{tabularx}
    \end{column}
    \begin{column}{0.42\textwidth}
      \begin{conceito}[Planos]
        \textbf{Spark}: gratuito, sem cartão; ao atingir a cota, o serviço para.\\
        \textbf{Blaze}: pague pelo uso, mantendo as cotas gratuitas.
      \end{conceito}
      \begin{atencao}[Cloud Storage exige Blaze]
        Desde \textbf{03/02/2026}, o Storage só funciona no plano Blaze. Configure um \textbf{alerta de orçamento}.
      \end{atencao}
      \begin{dica}
        O Módulo 16 compara o Firebase a outros BaaS.
      \end{dica}
    \end{column}
  \end{columns}
\end{frame}

% =============================================================================
\section{Configuração}

\begin{frame}[fragile]{FlutterFire CLI}
  \begin{columns}[T]
    \begin{column}{0.5\textwidth}
\begin{terminal}
$ npm install -g firebase-tools
$ firebase login
$ dart pub global activate flutterfire_cli
$ flutterfire configure
$ flutter pub add firebase_core firebase_auth \
    cloud_firestore
\end{terminal}
      \begin{conceito}[O que o \cod{configure} faz]
        Cria ou escolhe o projeto no console, registra os apps Android, iOS, web e macOS e gera \cod{lib/firebase\_options.dart} com as chaves de cada plataforma.
      \end{conceito}
    \end{column}
    \begin{column}{0.48\textwidth}
\begin{codigo}[title=main.dart]{Dart}
import 'package:firebase_core/firebase_core.dart';
import 'firebase_options.dart';

Future<void> main() async {
  WidgetsFlutterBinding.ensureInitialized();
  await Firebase.initializeApp(
    options: DefaultFirebaseOptions.currentPlatform,
  );
  runApp(const LabMonitorApp());
}
\end{codigo}
      \begin{dica}
        Rode \cod{flutterfire configure} de novo sempre que adicionar uma plataforma ou um plugin que exija configuração nativa.
      \end{dica}
    \end{column}
  \end{columns}
\end{frame}

% =============================================================================
\section{Authentication}

\begin{frame}[fragile]{E-mail e senha no repositório}
  \begin{columns}[T]
    \begin{column}{0.58\textwidth}
\begin{codigo}[title=data/repositories/auth\_repository.dart]{Dart}
class AuthRepository extends ChangeNotifier {
  AuthRepository() {
    _auth.authStateChanges().listen((_) => notifyListeners());
  }
  final _auth = FirebaseAuth.instance;
  User? get usuario => _auth.currentUser;
  bool get logado => usuario != null;

  Future<Result<void>> entrar(String email, String senha) async {
    try {
      await _auth.signInWithEmailAndPassword(
          email: email, password: senha);
      return const Ok(null);
    } on FirebaseAuthException catch (e) {
      return Falha(Exception(switch (e.code) {
        'invalid-credential' => 'E-mail ou senha incorretos.',
        _ => 'Falha no login (${e.code}).',
      }));
    }
  }
}
\end{codigo}
    \end{column}
    \begin{column}{0.4\textwidth}
      \begin{conceito}[Estado da sessão]
        \cod{authStateChanges()} emite o usuário ao entrar, ao sair e ao abrir o app (a sessão fica salva). O repositório vira a fonte da verdade do login.
      \end{conceito}
      \begin{dica}[Outras operações]
        \cod{createUserWithEmailAndPassword}, \cod{sendPasswordResetEmail} e \cod{signOut()}.
      \end{dica}
      \begin{atencao}
        Com a proteção contra enumeração de e-mails, usuário ou senha errados viram \cod{invalid-credential}.
      \end{atencao}
    \end{column}
  \end{columns}
\end{frame}

\begin{frame}[fragile]{Login com Google e proteção de rotas}
  \begin{columns}[T]
    \begin{column}{0.54\textwidth}
\begin{codigo}[title=Google (google\_sign\_in 7)]{Dart}
Future<UserCredential> entrarComGoogle() async {
  if (kIsWeb) {
    return FirebaseAuth.instance
        .signInWithPopup(GoogleAuthProvider());
  }
  await GoogleSignIn.instance.initialize();
  final conta = await GoogleSignIn.instance.authenticate();
  final cred = GoogleAuthProvider.credential(
      idToken: conta.authentication.idToken);
  return FirebaseAuth.instance.signInWithCredential(cred);
}
\end{codigo}
      \begin{atencao}
        Android: cadastre as impressões SHA-1 e SHA-256 da chave no console. iOS: adicione o esquema de URL do cliente no \cod{Info.plist}.
      \end{atencao}
    \end{column}
    \begin{column}{0.44\textwidth}
\begin{codigo}[title=go\_router (Módulo 07)]{Dart}
GoRouter(
  refreshListenable: auth,
  redirect: (ctx, state) {
    final noLogin =
        state.matchedLocation == '/login';
    if (!auth.logado) {
      return noLogin ? null : '/login';
    }
    return noLogin ? '/painel' : null;
  },
  routes: [ ... ],
);
\end{codigo}
      \begin{dica}[App Store]
        Quem oferece login social de terceiros precisa oferecer também uma opção equivalente focada em privacidade, como \textbf{Sign in with Apple}.
      \end{dica}
    \end{column}
  \end{columns}
\end{frame}

% =============================================================================
\section{Cloud Firestore}

\begin{frame}{Modelo de dados: coleções e documentos}
  \begin{columns}[c]
    \begin{column}{0.55\textwidth}
      \centering
      \begin{tikzpicture}[node distance=2.5mm]
        \node[caixapetroleo,minimum width=2.8cm] (c1) {\cod{laboratorios}};
        \node[caixa,below=of c1,minimum width=2.8cm,align=left,font=\scriptsize] (d1) {\textbf{lab-3}\\ nome: "Eletrônica"\\ bloco: "B"};
        \node[caixapetroleo,below left=4mm and -8mm of d1,minimum width=2.3cm] (c2) {\cod{bancadas}};
        \node[caixapetroleo,below right=4mm and -8mm of d1,minimum width=2.3cm] (c3) {\cod{leituras}};
        \node[caixa,below=of c2,minimum width=2.3cm,align=left,font=\scriptsize] (d2) {\textbf{b-07}\\ nome: "B7"\\ livre: false};
        \node[caixa,below=of c3,minimum width=2.3cm,align=left,font=\scriptsize] (d3) {\textbf{auto-id}\\ sensor: "temp"\\ valor: 23.5};
        \draw[seta] (c1) -- (d1); \draw[seta] (d1) -- (c2); \draw[seta] (d1) -- (c3);
        \draw[seta] (c2) -- (d2); \draw[seta] (c3) -- (d3);
      \end{tikzpicture}
    \end{column}
    \begin{column}{0.43\textwidth}
      \begin{conceito}[Estrutura]
        \textbf{Coleções} guardam \textbf{documentos}; documentos guardam campos e podem ter \textbf{subcoleções}. Cada documento tem até 1\,MiB.
      \end{conceito}
      \begin{dica}[Modele pelas telas]
        Sem \cod{JOIN}: guarde junto o que é lido junto. Duplicar um campo (como o nome do laboratório na reserva) é normal e intencional.
      \end{dica}
      \begin{atencao}
        Cobrança por leitura de documento: uma lista de 1\,000 itens custa 1\,000 leituras. Use \cod{limit} e paginação.
      \end{atencao}
    \end{column}
  \end{columns}
\end{frame}

\begin{frame}[fragile]{Lendo e gravando com \cod{withConverter}}
  \begin{columns}[T]
    \begin{column}{0.6\textwidth}
\begin{codigo}[title=data/repositories/bancada\_repository\_firestore.dart]{Dart}
class BancadaRepositoryFirestore {
  BancadaRepositoryFirestore(this.labId);
  final String labId;
  late final _col = FirebaseFirestore.instance
      .collection('laboratorios/$labId/bancadas')
      .withConverter<Bancada>(
        fromFirestore: (s, _) => Bancada.deMapa(s.id, s.data()!),
        toFirestore: (b, _) => b.paraMapa(),
      );

  Stream<List<Bancada>> observarLivres() => _col
      .where('livre', isEqualTo: true)
      .orderBy('nome').limit(50)
      .snapshots()
      .map((q) => q.docs.map((d) => d.data()).toList());

  Future<void> criar(Bancada b) => _col.add(b);
  Future<void> renomear(String id, String nome) =>
      _col.doc(id).update({'nome': nome});
  // remover: _col.doc(id).delete()
}
\end{codigo}
    \end{column}
    \begin{column}{0.38\textwidth}
      \begin{conceito}[\cod{snapshots()}]
        \cod{Stream} que emite a lista sempre que um documento muda, em qualquer aparelho.
      \end{conceito}
      \begin{conceito}[Offline]
        Android e iOS mantêm cache local: lê sem rede e envia as gravações quando ela volta.
      \end{conceito}
      \begin{dica}
        Filtro e ordenação em campos diferentes pedem um \textbf{índice composto}.
      \end{dica}
    \end{column}
  \end{columns}
\end{frame}

\begin{frame}[fragile]{Transações, lotes e paginação}
  \begin{columns}[T]
    \begin{column}{0.56\textwidth}
\begin{codigo}[title=Reserva sem conflito]{Dart}
Future<void> reservar(String id, String uid) =>
    FirebaseFirestore.instance.runTransaction((tx) async {
      final ref = _col.doc(id);
      final snap = await tx.get(ref);
      if (snap.data()?.livre != true) {
        throw StateError('Bancada já reservada');
      }
      tx.update(ref, {
        'livre': false,
        'reservadaPor': uid,
        'reservadaEm': FieldValue.serverTimestamp(),
      });
    });
\end{codigo}
\begin{codigo}[title=Próxima página]{Dart}
final pagina = await _col.orderBy('nome')
    .startAfterDocument(ultimo).limit(20).get();
\end{codigo}
    \end{column}
    \begin{column}{0.42\textwidth}
      \begin{conceito}[Transação]
        Lê e grava de forma atômica. Se outro aparelho alterar a bancada no meio, o Firestore repete a função. Dois alunos nunca reservam a mesma bancada.
      \end{conceito}
      \begin{conceito}[Lote (\cod{batch})]
        Várias gravações tudo-ou-nada, sem leitura: \cod{batch()}, \cod{set}, \cod{update}, \cod{delete} e \cod{commit()}.
      \end{conceito}
      \begin{dica}
        \cod{FieldValue.serverTimestamp()} usa o relógio do servidor, não o do celular.
      \end{dica}
    \end{column}
  \end{columns}
\end{frame}

% =============================================================================
\section{Regras de segurança}

\begin{frame}[fragile]{Regras do Firestore}
  \begin{columns}[T]
    \begin{column}{0.58\textwidth}
\begin{codigo}[title=firestore.rules]{JSON}
rules_version = '2';
service cloud.firestore {
  match /databases/{database}/documents {
    function logado() { return request.auth != null; }

    match /laboratorios/{lab}/bancadas/{id} {
      allow read: if logado();
      allow update: if logado() && request.resource.data
        .diff(resource.data).affectedKeys()
        .hasOnly(['livre', 'reservadaPor', 'reservadaEm']);
      allow create, delete:
        if request.auth.token.admin == true;
    }
    match /laboratorios/{lab}/leituras/{id} {
      allow read: if logado();
      allow create: if logado()
        && request.resource.data.valor is number;
    }
  }
}
\end{codigo}
    \end{column}
    \begin{column}{0.4\textwidth}
      \begin{erro}[O erro mais grave]
        Publicar com \cod{allow read, write: if true;} (``modo de teste''). Qualquer pessoa lê e apaga todo o banco.
      \end{erro}
      \begin{conceito}[Como pensar]
        O app \textbf{não é confiável}: qualquer um pode chamar a API com o seu \cod{firebase\_options}. As regras, no servidor, são a barreira real.
      \end{conceito}
      \begin{dica}
        \cod{admin} é uma \emph{custom claim} definida por uma Cloud Function ou pelo Admin SDK.
      \end{dica}
    \end{column}
  \end{columns}
\end{frame}

% =============================================================================
\section{Storage, Functions e mensagens}

\begin{frame}[fragile]{Storage e Cloud Functions}
  \begin{columns}[T]
    \begin{column}{0.52\textwidth}
\begin{codigo}[title=Foto do equipamento]{Dart}
Future<String> enviarFoto(String id, File foto) async {
  final ref = FirebaseStorage.instance
      .ref('equipamentos/$id.jpg');
  final tarefa = ref.putFile(foto,
      SettableMetadata(contentType: 'image/jpeg'));
  tarefa.snapshotEvents.listen((s) => progresso.value =
      s.bytesTransferred / s.totalBytes);
  await tarefa;
  return ref.getDownloadURL();
}
\end{codigo}
      \begin{dica}
        O Storage tem regras próprias (\cod{storage.rules}): limite tamanho e tipo do arquivo.
      \end{dica}
    \end{column}
    \begin{column}{0.46\textwidth}
\begin{codigo}[title=Chamando uma função]{Dart}
final gerar = FirebaseFunctions.instance
    .httpsCallable('gerarRelatorio');
final r = await gerar.call({'lab': 'lab-3',
    'mes': '2026-10'});
final url = r.data['url'] as String;
\end{codigo}
      \begin{conceito}[Cloud Functions]
        Código em Node.js ou Python que roda no servidor: lógica que não pode ficar no app, envio de push, integração com outros serviços. Exige plano Blaze.
      \end{conceito}
    \end{column}
  \end{columns}
\end{frame}

\begin{frame}[fragile]{Notificações push com FCM}
  \begin{columns}[T]
    \begin{column}{0.58\textwidth}
\begin{codigo}[title=Configurando o Firebase Cloud Messaging]{Dart}
@pragma('vm:entry-point')
Future<void> _emSegundoPlano(RemoteMessage m) async {
  await Firebase.initializeApp(
      options: DefaultFirebaseOptions.currentPlatform);
  // gravar no banco local, atualizar contador...
}

Future<void> configurarPush() async {
  FirebaseMessaging.onBackgroundMessage(_emSegundoPlano);
  final fcm = FirebaseMessaging.instance;
  await fcm.requestPermission();        // iOS e Android 13+
  final token = await fcm.getToken();   // enviar ao backend
  fcm.onTokenRefresh.listen(salvarToken);
  FirebaseMessaging.onMessage.listen(mostrarLocal);
  FirebaseMessaging.onMessageOpenedApp.listen(
      (m) => router.go('/alertas/${m.data['id']}'));
}
\end{codigo}
    \end{column}
    \begin{column}{0.4\textwidth}
      \begin{conceito}[Três situações]
        \textbf{Primeiro plano}: \cod{onMessage}; você exibe.\\
        \textbf{Fundo ou fechado}: o sistema exibe; o \emph{handler} roda isolado.\\
        \textbf{Toque}: \cod{onMessageOpenedApp} abre a tela.
      \end{conceito}
      \begin{atencao}
        iOS exige a chave APNs no console e a capacidade \emph{Push Notifications} no Xcode.
      \end{atencao}
    \end{column}
  \end{columns}
\end{frame}

% =============================================================================
\section{Monitoramento e desenvolvimento}

\begin{frame}[fragile]{Crashlytics, Analytics e Remote Config}
  \begin{columns}[T]
    \begin{column}{0.56\textwidth}
\begin{codigo}[title=main.dart: capturando falhas]{Dart}
FlutterError.onError =
    FirebaseCrashlytics.instance.recordFlutterFatalError;

PlatformDispatcher.instance.onError = (erro, pilha) {
  FirebaseCrashlytics.instance
      .recordError(erro, pilha, fatal: true);
  return true;
};
\end{codigo}
\begin{codigo}[title=Evento e configuração remota]{Dart}
await FirebaseAnalytics.instance.logEvent(
    name: 'reserva_criada', parameters: {'lab': 'lab-3'});

final rc = FirebaseRemoteConfig.instance;
await rc.fetchAndActivate();
final intervalo = rc.getInt('intervalo_leitura_s');
\end{codigo}
    \end{column}
    \begin{column}{0.42\textwidth}
      \begin{conceito}[Para quê]
        \textbf{Crashlytics}: falhas com pilha e dispositivo.\\
        \textbf{Analytics}: o que os usuários fazem.\\
        \textbf{Remote Config}: mudar valores e ligar recursos sem nova versão na loja.\\
        \textbf{App Check}: recusa chamadas que não vêm do seu app.
      \end{conceito}
      \begin{atencao}[LGPD]
        Informe na política de privacidade os dados coletados e não envie dados pessoais em eventos.
      \end{atencao}
    \end{column}
  \end{columns}
\end{frame}

\begin{frame}[fragile]{Emulator Suite}
  \begin{columns}[T]
    \begin{column}{0.5\textwidth}
\begin{terminal}
$ firebase init emulators
$ firebase emulators:start
  Auth      localhost:9099
  Firestore localhost:8080
  Storage   localhost:9199
  UI        localhost:4000
\end{terminal}
\begin{codigo}[title=Apontando o app para o emulador]{Dart}
if (kDebugMode) {
  final host = defaultTargetPlatform ==
      TargetPlatform.android ? '10.0.2.2' : 'localhost';
  await FirebaseAuth.instance.useAuthEmulator(host, 9099);
  FirebaseFirestore.instance
      .useFirestoreEmulator(host, 8080);
}
\end{codigo}
    \end{column}
    \begin{column}{0.48\textwidth}
      \begin{conceito}[Vantagens]
        Sem custo, sem tocar nos dados reais; testa as regras antes de publicar e serve aos testes de integração (Módulo 14).
      \end{conceito}
      \begin{dica}
        No emulador Android, \cod{10.0.2.2} é o \cod{localhost} do computador.
      \end{dica}
      \begin{atencao}
        Não use o Firebase direto nos widgets: encapsule em repositórios (Módulo 08).
      \end{atencao}
    \end{column}
  \end{columns}
\end{frame}

% =============================================================================
\section{Fechamento}

\begin{frame}{Resumo}
  \begin{resumo}
    \begin{itemize}
      \item \cod{flutterfire configure} gera as opções; \cod{Firebase.initializeApp} inicia tudo.
      \item Authentication com e-mail, Google e Apple; \cod{authStateChanges} protege as rotas.
      \item Firestore: coleções, documentos, \cod{snapshots()} em tempo real, transações e cache offline.
      \item Regras de segurança são a única proteção real dos dados.
      \item Storage (Blaze), Functions, FCM, Crashlytics e Remote Config completam o backend.
    \end{itemize}
  \end{resumo}
\end{frame}

\begin{frame}{Glossário (1/2)}
  \begin{glossario}
    BaaS & \emph{backend} pronto oferecido como serviço: contas, banco, arquivos \\
    FlutterFire & plugins oficiais do Firebase para Flutter \\
    OAuth 2.0 & delegação de autorização por \emph{tokens}, sem expor a senha \\
    OpenID Connect & camada de identidade sobre o OAuth (ID \emph{token}) \\
    PKCE & segredo por login que impede o uso de um código interceptado \\
    ID \emph{token} & JWT que diz quem é o usuário; vale 1 hora no Firebase \\
    \emph{Custom claims} & papéis gravados no \emph{token} pelo servidor \\
    \emph{Passkey} & credencial de par de chaves, resistente a \emph{phishing} \\
    \cod{flutterfire configure} & gera o \cod{firebase\_options.dart} de cada plataforma \\
  \end{glossario}
\end{frame}

\begin{frame}{Glossário (2/2)}
  \begin{glossario}
    Cloud Firestore & banco de documentos com sincronização em tempo real \\
    \cod{snapshots()} & fluxo com o resultado de uma consulta e cada mudança \\
    Compensação de latência & a gravação aparece antes da confirmação do servidor \\
    Regras de segurança & autorização declarativa avaliada no servidor a cada requisição \\
    FCM & serviço de mensagens \emph{push} do Firebase \\
    \emph{Token} de registro & endereço do app em um aparelho, para receber \emph{push} \\
    Idempotência & processar o mesmo evento duas vezes dá o mesmo resultado \\
    Emulator Suite & emuladores locais de Auth, Firestore, Functions e Storage \\
  \end{glossario}
\end{frame}

\begin{frame}{Questões de revisão}
  \begin{questoes}
    \item Diferencie OAuth 2.0 de OpenID Connect. Qual deles o ``Entrar com Google'' usa?
    \item Por que apps nativos devem fazer login no navegador do sistema, com PKCE?
    \item Para que servem o ID \emph{token}, o \emph{token} de acesso e o \emph{token} de renovação?
    \item Por que \emph{passkeys} resistem a \emph{phishing} e senhas não?
    \item O que é compensação de latência? Como a tela sabe que um documento ainda está pendente?
    \item Por que as regras do Firestore são a verdadeira proteção dos dados? O que significa ``regras não são filtros''?
    \item Descreva o caminho de um alerta \emph{push} da função até um iPhone e até um Android.
    \item Por que uma função acionada por evento precisa ser idempotente?
  \end{questoes}
\end{frame}

\begin{frame}{Exercícios}
  \begin{exercicio}[Exercício 1 --- Autenticação]
    Implemente cadastro, login, recuperação de senha e saída, com mensagens de erro em português.
  \end{exercicio}
  \begin{exercicio}[Exercício 2 --- Regras]
    Escreva regras em que só o autor pode cancelar a própria reserva e teste-as no Emulator Suite.
  \end{exercicio}
  \begin{exercicio}[Exercício 3 --- Push]
    Envie uma notificação pelo console e faça o toque abrir a tela de alertas.
  \end{exercicio}
  \begin{projeto}[LabMonitor --- Etapa 12]
    Login com e-mail e Google, bancadas e reservas no Firestore com transação, fotos no Storage e Crashlytics ativo.
  \end{projeto}
\end{frame}

\begin{frame}{Para aprofundar}
  \small
  \begin{itemize}
    \item FlutterFire: \url{https://firebase.flutter.dev}
    \item Firebase para Flutter: \url{https://firebase.google.com/docs/flutter/setup}
    \item Authentication: \url{https://firebase.google.com/docs/auth/flutter/start}
    \item Cloud Firestore: \url{https://firebase.google.com/docs/firestore}
    \item Regras de segurança: \url{https://firebase.google.com/docs/rules}
    \item Emulator Suite: \url{https://firebase.google.com/docs/emulator-suite}
  \end{itemize}
\end{frame}

\begin{frame}{Respostas das questões de revisão}
  \begin{respostas}
    \item O OAuth 2.0 delega autorização (\emph{token} de acesso); o OpenID Connect acrescenta identidade (ID \emph{token}). O login com Google usa OpenID Connect: o app recebe um ID \emph{token} e o troca por uma credencial do Firebase.
    \item O navegador do sistema mantém a senha longe do app e reaproveita a sessão já aberta; o PKCE impede que um código interceptado seja trocado por \emph{tokens}.
    \item ID \emph{token}: quem é o usuário. Acesso: autoriza chamadas a uma API. Renovação: obtém novos \emph{tokens} sem novo login.
    \item A \emph{passkey} só funciona no domínio para o qual foi criada e assina um desafio sem enviar segredo. A senha pode ser digitada num site falso e reutilizada.
    \item A gravação aparece na hora, antes da confirmação do servidor; se ele recusar, o SDK desfaz a mudança. Enquanto isso, \cod{metadata.hasPendingWrites} é verdadeiro.
    \item O app fala direto com o banco e qualquer um pode usar as chaves públicas: só as regras, avaliadas no servidor, barram o acesso. Uma consulta que \textbf{possa} trazer documentos proibidos falha inteira; ela precisa ter os mesmos filtros da regra.
    \item A função chama o FCM pela API HTTP v1 com o \emph{token} do aparelho. No Android, o FCM entrega pelo Google Play services; no iPhone, repassa ao APNs, que entrega pela conexão do sistema.
    \item O provedor garante pelo menos uma execução por evento e às vezes repete; sem idempotência, o mesmo alerta seria enviado duas vezes.
  \end{respostas}
\end{frame}

\encerramento{Módulo 13 --- Recursos do dispositivo: câmera, GPS, sensores, Bluetooth, NFC, notificações e IoT.}

\end{document}
